\clearpage\pagestyle{otMain}\fontsize{11}{15}\selectfont
\corrigesChapitre{16}{Raisonner et faire la synthèse}
\begin{blocLivre}{cor-16-1}
\begin{corrige}{1}{Appartenance et inclusion}
\demarche{Distinguer un élément de l’ensemble qui le contient, puis appliquer chaque opération sur les ensembles.}
Les deux premières phrases sont vraies : 2 est un élément de $A$, et le seul élément de $\{2\}$ est dans $A$. La troisième est fausse car 3 n’est pas listé. La dernière est vraie : aucun élément de l’ensemble vide ne peut contredire l’inclusion. Les symboles $\in$ et $\subset$ ne décrivent pas la même relation. Ensuite $\res{A\cap B=\{2\}}$, $\res{A\cup B=E}$ et $\res{\overline A=\{3;5\}}$. Donc $\res{\operatorname{Card}(A\cup B)=5}$, ce que confirme $3+3-1=5$.
\end{corrige}
\end{blocLivre}
\begin{blocLivre}{cor-16-2}
\begin{corrige}{2}{Un produit cartésien}
\demarche{Associer chaque élément du premier ensemble à tous ceux du second.}
On obtient $(0;a),(0;b),(0;c),(1;a),(1;b),(1;c)$. Il y a $\res6$ couples, soit $2\times3$. L’ordre des composantes compte : les premières appartiennent à $A$ et les secondes à $B$. Le produit cartésien est un outil naturel pour décrire deux épreuves successives.
\end{corrige}
\end{blocLivre}
\begin{blocLivre}{cor-16-3}
\begin{corrige}{3}{Nier une condition}
\demarche{Compléter toutes les possibilités sans laisser de borne oubliée.}
La négation de $x<5$ est $\res{x\ge5}$, en incluant le cas d’égalité. La négation des deux conditions réunies est « $\res{x<2\text{ ou }x>7}$ ». Il suffit que l’une des deux contraintes échoue pour sortir de l’intervalle $[2;7]$.
\end{corrige}
\end{blocLivre}
\begin{blocLivre}{cor-16-4}
\begin{corrige}{4}{Implication et réciproque}
\demarche{Prouver la phrase initiale puis tester sa réciproque indépendamment.}
L’implication est vraie : $n=6k=2(3k)$ est pair. La réciproque est « si un entier est pair, il est multiple de 6 ». Elle est fausse, par exemple pour 2. Une hypothèse plus forte peut entraîner une propriété sans que cette propriété suffise à retrouver l’hypothèse.
\end{corrige}
\end{blocLivre}
\begin{blocLivre}{cor-16-5}
\begin{corrige}{5}{Prouver une existence}
\demarche{Identifier ce que demande la formulation avant de choisir une preuve.}
On résout $x(x-2)=0$ : les solutions sont $\res{0\text{ et }2}$. La phrase d’existence est donc vraie. La phrase « pour tout réel $x$, $x^2=2x$ » est fausse, par exemple en $x=1$. Un seul exemple suffit pour établir une existence, mais il ne démontre pas une propriété universelle.
\end{corrige}
\end{blocLivre}
\begin{blocLivre}{cor-16-6}
\begin{corrige}{6}{Un contre-exemple décisif}
\demarche{Chercher un oubli de signe dans une affirmation sur les carrés.}
Prendre $x=-3$ : son carré vaut 9, supérieur à 4, mais $-3$ n’est pas supérieur à 2. La conclusion correcte est $\res{x<-2\text{ ou }x>2}$. Elle tient compte des deux branches de la fonction carré. On peut aussi formuler la conclusion comme une distance à 0 supérieure à 2.
\end{corrige}
\end{blocLivre}
\begin{blocLivre}{cor-16-7}
\begin{corrige}{7}{Une disjonction de cas}
\demarche{Faire deux preuves partielles qui couvrent toutes les possibilités.}
Si $n=2k$, alors $n(n+1)=2k(n+1)$ est pair. Si $n=2k+1$, alors $n+1=2(k+1)$ et $n(n+1)=2n(k+1)$ est pair. Les deux cas couvrent tous les entiers, donc la propriété est établie pour tout $n\in\mathbb Z$. Il faut montrer que la liste des cas est exhaustive.
\end{corrige}
\end{blocLivre}
\begin{blocLivre}{cor-16-8}
\begin{corrige}{8}{Une preuve par l’absurde}
\demarche{Supposer l’existence de l’objet impossible puis construire une contradiction.}
Supposons qu’un plus grand entier naturel $N$ existe. Le nombre $N+1$ est encore un entier naturel, et il est strictement supérieur à $N$. Cela contredit le choix de $N$. L’hypothèse de départ est donc impossible. Le raisonnement ne dépend d’aucune valeur particulière de $N$.
\end{corrige}
\end{blocLivre}
\begin{blocLivre}{cor-16-9}
\begin{corrige}{9}{Une contraposée}
\demarche{Écrire correctement la contraposée avant le calcul.}
La contraposée est « si $n$ est pair, alors $n^2$ est pair ». Écrire $n=2k$ donne $n^2=4k^2=2(2k^2)$, donc un entier pair. La contraposée étant démontrée, l’implication initiale est vraie. Cette méthode évite de supposer directement l’écriture recherchée de $n$.
\end{corrige}
\end{blocLivre}
\begin{blocLivre}{cor-16-10}
\begin{corrige}{10}{Le statut d’un dessin}
\demarche{Distinguer une observation numérique d’une condition mathématique exacte.}
Il faut par exemple des coordonnées exactes et un déterminant nul, ou une propriété géométrique démontrée. Les points $(0;0),(1;1),(2;2{,}001)$ paraissent presque alignés à petite échelle, mais le déterminant vaut $1\times2{,}001-1\times2=0{,}001\ne0$. Ils ne sont donc pas alignés. La précision d’affichage ne remplace pas l’égalité exacte.
\end{corrige}
\end{blocLivre}
\begin{blocLivre}{cor-16-11}
\begin{corrige}{11}{Synthèse : une aire maximale}
\demarche{Relier une contrainte géométrique à une optimisation algébrique.}
La somme des deux côtés vaut 14 ; l’autre côté est $14-x$ avec $0<x<14$. L’aire est $A(x)=x(14-x)=49-(x-7)^2\le49$. Le maximum $\res{49\text{ cm}^2}$ est atteint pour $\res{7\text{ cm}\times7\text{ cm}}$. Le rectangle optimal est un carré ; les dimensions respectent le périmètre annoncé.
\label{fig-cor-16-11}
\figureGeo[corrige]{synthese-rectangle}
\end{corrige}
\end{blocLivre}
\begin{blocLivre}{cor-16-12}
\begin{corrige}{12}{Synthèse : contrôler une pièce}
\demarche{Construire les effectifs avant d’inverser le conditionnement.}
La sensibilité vaut $36/40=\res{90\,\%}$. Il y a $36+23=59$ signaux positifs, dont 36 vrais défauts ; la probabilité inversée est $\res{36/59\approx61{,}0\,\%}$. La fréquence initiale des défauts est $40/500=8\,\%$, distincte des deux autres valeurs. Chaque pourcentage répond à une question de référence différente.
\end{corrige}
\end{blocLivre}
\begin{blocLivre}{cor-16-13}
\begin{corrige}{13}{Synthèse : coûts et seuil}
\demarche{Comparer recettes et dépenses puis interpréter une inégalité entière.}
Les recettes sont $18n$, les dépenses $180+12n$ ; le résultat est $B(n)=6n-180$. Il est positif ou nul pour $n\ge30$. Le seuil est donc $\res{30\text{ participants}}$, avec résultat nul au seuil ; un bénéfice strict demande au moins 31 participants. Le domaine est l’ensemble des entiers naturels.
\end{corrige}
\end{blocLivre}
\begin{blocLivre}{cor-16-14}
\begin{corrige}{14}{Synthèse : un lieu et une droite}
\demarche{Choisir d’abord la propriété géométrique qui simplifie le système.}
Les points équidistants de $A$ et $B$ sont sur la médiatrice $x=2$. L’intersection avec $y=x+1$ est donc $\res{M(2;3)}$. Contrôle : $MA^2=4+9=13$ et $MB^2=4+9=13$. Le problème se résout plus vite en utilisant la médiatrice qu’en développant d’emblée toutes les distances.
\label{fig-cor-16-14}
\figureGeo[corrige]{synthese-mediatrice}
\end{corrige}
\end{blocLivre}
\begin{blocLivre}{cor-16-15}
\begin{corrige}{15}{Synthèse : juger une conclusion}
\demarche{Construire un exemple qui respecte l’hypothèse mais contredit la conclusion.}
Prendre une mesure égale à $-9$ et 19 mesures égales à 11. La somme vaut $-9+19\times11=200$, donc la moyenne vaut bien 10. Pourtant une seule mesure est inférieure ou égale à 10. L’affirmation est fausse. C’est la médiane, et non la moyenne, qui décrit un partage de la série ordonnée en deux parties.
\end{corrige}
\end{blocLivre}
